% ID: SEL-D0056-OB-E02
% ESTADO: validada. S057.
\textbf{a)} El dominio es $\mathbb R\setminus\{-2,2\}$. Cada rama racional es continua y derivable donde está definida, salvo que en el punto de unión hay que comparar los límites. En $0$ ambos límites y $f(0)$ valen $-1$, así que hay continuidad. Las derivadas de las ramas son $2/(x+2)^2$ para $x<0$ y $-2/(x-2)^2$ para $x>0$, excluyendo los polos. En $0$ las derivadas laterales son $1/2$ y $-1/2$, por lo que $f$ es derivable en su dominio excepto en $0$.

\textbf{b)} Los límites en $-2$ son $+\infty$ por la izquierda y $-\infty$ por la derecha. En $2$ son $-\infty$ por la izquierda y $+\infty$ por la derecha. Las asíntotas verticales son $x=-2$ y $x=2$. Como los límites en ambos infinitos son cero, la asíntota horizontal es $y=0$.
